\section{Synthesis: what carries forward}
\label{m:sec:synthesis}

Four things have been assembled in this chapter, and it is worth separating them
before they are used.

The first is standard apparatus: exponential clocks and the races between them,
continuous-time Markov chains and their generators, first-step analysis for hitting
probabilities, generating functions, and the linear birth--death process together
with the first-order equation \cref{m:eq:bdpde} satisfied by its generating
function. That equation is the template for every generating-function calculation
in this thesis.

The second is a working method rather than a result, and it is the one used most
often. A model is specified by rates; the rates give a generator and a master
equation; multiplying by $z^i$, or by $x^iy^j$, and summing converts that master
equation into a single first-order partial differential equation for a generating
function; and the method of characteristics of \cref{m:sec:mocabsdeath} solves
equations of that type. \Cref{m:eq:Gabd} shows the route completed on a problem for
which no shortcut exists, and the same route is followed
\fwd{bdc}{in the later chapters} when a catastrophe channel is added to the rates
of \cref{m:eq:bdrates}. \Cref{m:app:extract} supplies the last step, recovering
time-dependent state probabilities from a generating function given by a formula
rather than a series.

The third is the quasi-stationary picture. A subcritical binary Galton--Watson
process is certain to die out, but conditioned on survival its population settles
to a fixed law whose mean is $1/A(p)$ and whose variance is \cref{m:eq:condvar}.
Both are expressed entirely in terms of $p$ and the constant
$A(p) = \lim_n S_n/(2p)^n$. The continuous-time analogue $\Ac(p) = (1-2p)/(1-p)$ is
elementary, and \cref{m:sec:cts} exhibits it in full, along with the structural
reason the discrete constant should be harder: in continuous time the survival
probability is elementary at every time, whereas in discrete time the amplitude is
the accumulated product of every early nonlinear correction. Everything further
about $A(p)$ --- its exact product and series representations, the two-sided bounds
that trap it between $\tfrac12\Ac(p)$ and $2\varepsilon/(1+2\varepsilon)$, the
near-critical asymptotic $A(p)\sim2(1-2p)$, the failed search for a closed form, and
the identification of $A(p)$ with a value of the Koenigs linearising function of the
logistic map --- belongs to \ChA{} and is developed there.

The fourth is a definition and a framework held in reserve. The
birth--death--catastrophe process of \cref{m:def:bdc} has been defined and its
generating-function equation \cref{m:eq:bdcpde} written down, and no more; the
killed-generator viewpoint of \cref{m:sec:killedchain} places it and its relatives
in a common setting, with the quasi-stationary law appearing as a left eigenvector
of the killed generator and the survival decay rate as its eigenvalue. The results
themselves belong to the chapters that need them.

\subsection*{Questions left open here}

Three questions raised by the material above are not settled in this chapter, and
two of them are proper to it.

Whether a quasi-stationary distribution survives when a birth--death--catastrophe
process is conditioned on \emph{neither} extinction nor rupture is not settled by
the argument of \cref{m:sec:bdc}. The two absorbing mechanisms are not
interchangeable --- \cref{m:eq:bdcsurvivalloss} shows that ordinary death drains
survival mass only from state $1$ while catastrophe drains it from every positive
state --- and the joint conditioning is best approached by simulation in the first
instance.

How tight the Jensen bound \cref{m:eq:bdcguess} is remains a numerical question. A
sharp answer would make the discrete catastrophe calculation usable in its own
right rather than merely indicative; \cref{m:rem:notmeanfield} explains why no
mean-field substitution will supply one.

The third question belongs to \ChA{} and is recorded here only so that it is not
looked for in this chapter: whether the map $p\mapsto A(p)$, as opposed to the
Koenigs function in its own argument, admits an elementary description at all.

\FloatBarrier
